\documentclass[10pt,a4paper,UTF8,fontset=fandol]{ctexart}
\usepackage[a4paper,left=1.18cm,right=1.18cm,top=1.05cm,bottom=1.05cm]{geometry}
\usepackage{tabularx}
\usepackage{xcolor}
\usepackage[hidelinks]{hyperref}
\usepackage{titlesec}
\usepackage{enumitem}
\usepackage{microtype}

\definecolor{navy}{HTML}{173B63}
\definecolor{ruleblue}{HTML}{4E79A7}
\definecolor{muted}{HTML}{4B5968}
\setlength{\parindent}{0pt}
\setlength{\parskip}{0pt}
\pagestyle{empty}
\linespread{1.10}
\setlist[itemize]{leftmargin=1.25em,itemsep=0.30em,topsep=0.18em,parsep=0pt,partopsep=0pt}
\titleformat{\section}{\fontsize{13.0}{15.0}\selectfont\bfseries\color{navy}}{}{0pt}{}[\vspace{-0.34em}\color{ruleblue}\rule{\linewidth}{0.65pt}]
\titlespacing*{\section}{0pt}{0.85em}{0.34em}

\newcommand{\entry}[3]{%
  \begin{tabularx}{\linewidth}{@{}X r@{}}
    \textbf{\color{navy}#1} & {\color{muted}#2} \\
    \multicolumn{2}{@{}l@{}}{{\fontsize{8.9}{10.4}\selectfont\color{muted}#3}}
  \end{tabularx}\par\vspace{-0.06em}
}


\begin{document}
\fontsize{9.5}{12.0}\selectfont
{\fontsize{23.5}{26}\selectfont\bfseries 潘佳祥}\hfill {\color{navy}\bfseries 平台可靠性 Infra 实习生}
\par\vspace{0.3em}{\color{muted}上海｜13501632852｜pym66@outlook.com｜github.com/pym96}
\par\vspace{0.6em}{\color{navy}\bfseries 候选草稿·禁止投递}
\par\vspace{0.3em}以任务状态、失败恢复和受控负载为重点的候选定位。
\section{项目经历}
{\bfseries\color{navy}Pi Agent Platform｜私有候选项目材料}\par
{\color{muted}\fontsize{8.8}{10.8}\selectfont Go / PostgreSQL / Pi / Docker Compose｜Agent协作实施的服务实验，非本人独立实现。}\par
\begin{itemize}
% P01a
\item PG保存任务与有序事件，幂等键区分重复提交；未知执行阻断会话，不自动重发模型请求。
% P01b
\item 终态与事件同事务提交；单控制实例持锁，两个worker执行，事件读取与控制连接池分离。
% P01c
\item Compose分离API、PG、Pi与负载端，合计4 CPU / 8 GiB硬配额；数据卷独立、端口仅loopback。
% P02a
\item hey固定并发、Vegeta固定目标速率：54轮共72600次请求均HTTP 200，分开核实际发送与完成率。
% P02b
\item SSE在120s内24次采样均为10000在线订阅；断连后API观测回到0连接、11条goroutine。
\end{itemize}
{\color{muted}\fontsize{8.5}{10.4}\selectfont 测量边界：synthetic端点、单控制实例、最多两个执行者；HTTP history仅两行，非模型吞吐。\par}
{\color{muted}\fontsize{8.5}{10.4}\selectfont SSE仅一次120s采样逻辑订阅；含1000次计划重连、21次服务写失败，不代表TCP不中断或生产容量。\par}
\vspace{0.6em}
{\bfseries\color{navy}Workspace Agent Harness}\par
{\color{muted}\fontsize{8.8}{10.8}\selectfont 面向不确定Agent执行的Runtime、协议适配与可重算实验｜个人公开项目｜2026.08--至今}\par
\begin{itemize}
% A01
\item 权限取Runtime、caller与Pack交集；evaluator在有界独立进程运行，错误不改写执行终态。
% A02
\item 离线typed-history Adapter校验关联；事件化AgentLoop先验证动作，再写有序事件与唯一终态。
% A05
\item TUI从保留事件零模型/工具调用回放，取消只记一个终态；c2c9e33基线仓库测试116/116。
\end{itemize}
\vspace{0.6em}
{\bfseries\color{navy}RoboMaster}\par
{\color{muted}\fontsize{8.8}{10.8}\selectfont RoboMaster自动瞄准与自主哨兵｜主要提交者｜2023.06--2024.07}\par
\begin{itemize}
% R01
\item Git identity覆盖MA\_Sentry的179/188与MA\_training的144/144提交；包含PnP、串口与ROS工程。
\end{itemize}
\vspace{0.6em}
\section{教育背景}
上海大学｜控制科学与工程｜硕士｜2025.09--2028.06｜GPA 3.76/4.00
\section{技能}
编程与工程：Python、C++、ROS1、Git、unittest、JSON/JSONL\par
Agent：Agent Runtime、Tool Calling、JSON Schema、Trace、Evaluation\par
\section{公开输出}
GitHub 个人非 Fork 仓库累计 100+ stars。
\vfill{\color{muted}\fontsize{8.5}{10}\selectfont 新增Pi项目句子待本人披露批准；本稿仅供内部评审。}
\end{document}
